% /root/monodromy, 19 November 2026.
% Independently audited in full by /root, including the determinant seeds.
\subsection{Hypersurfaces in projective space in odd characteristic}
\label{subsec:projective-space-odd-monodromy}
\status{Proved; independently audited by /root, including both determinant seeds.}

Fix an odd prime $p$ and an algebraically closed field $k$ of
characteristic $p$.  For degree $d$ hypersurfaces in $\mathbf P^n$,
let $W_{n,d}$ be the Hasse--Witt invertibility open in the full
projective space of equations, and let
\[
 h_{n,d}=\binom{d-1}{n},\qquad
 V_{n,d}=\left(H^{n-1}(\mathcal O_Y),F=1\right).
\]
The binomial coefficient is zero for $d\leq n$.
The full coherent bundle has this constant rank for all nonzero
equations, including singular equations.  A general-position union
of $d$ hyperplanes belongs to $W_{n,d}$, so this open is nonempty.
The desired smooth ordinary locus is a dense open in $W_{n,d}$;
passing to it does not change the finite monodromy image.
All the comparison families below therefore take place inside
one explicit constant-rank Lang local system, except for a
valuation-one Hasse degeneration used solely to compute its tame
determinant inertia.

\subsubsection{A connectedness lemma for linear projection}

\begin{lemma}\label{lem:linear-projection-monodromy}
Let $U$ be a nonempty open in $\mathbf P^N_k$, and let
$T\to U$ be a connected finite etale cover.  Let
$f:\mathbf P^N\dashrightarrow\mathbf P^M$ be a linear projection
whose center $Z$ meets $U$.  The restriction of $T$ to the
geometric generic fibre of $f|_{U\setminus Z}$ is connected.
Consequently a finite local system on $U$ retains its full
monodromy after restriction to that geometric generic fibre.
\end{lemma}

\begin{proof}
The open $U$ is smooth and integral, so $T$ is smooth and integral.
Let $B=\operatorname{Bl}_Z(\mathbf P^N)$ and put
$T'=T\times_U B_U$, where $B_U$ is the inverse image of $U$.
This is the blowup of $T$ in the smooth etale inverse image of
$Z\cap U$; in particular it is smooth and integral.  The resolved
projection gives a morphism $T'\to\mathbf P^M$.
Choose $z\in Z\cap U$.  The exceptional fibre above $z$ is
$E_z\simeq\mathbf P^M$, mapping isomorphically to the target.
Its pullback to $T'$ is the disjoint union of copies of $E_z$,
since it is the product of $E_z$ with the finite reduced set $T_z$.

Put $K=k(\mathbf P^M)$, and let $L$ be the algebraic closure of $K$
inside $k(T')$.  It is a finite extension of $K$.  Let $C$ be the
normalization of $\mathbf P^M$ in $L$; it is finite and integral.
There is a morphism $T'\to C$ over $\mathbf P^M$.  Indeed, on an
affine open of the target, each element of its integral closure
in $L$ is integral over the local ring at every point of the
normal scheme $T'$ lying above this open.  It therefore belongs
to each such local ring and defines a regular function there.
These morphisms glue.  One copy of $E_z$ in $T'$ consequently
gives a section $\mathbf P^M\to C$ of the finite map.
Its image is closed, has dimension $M$, and lies in the integral
scheme $C$ of dimension $M$; it is all of $C$.  Thus $C\simeq
\mathbf P^M$ and $L=K$.

On the nonempty open $T|_{U\setminus Z}$, the map to $\mathbf P^M$
is a composition of a finite etale map and a smooth linear
projection.  Hence $k(T')/K$ is separably generated.  The equality
$L=K$ implies that this extension is regular, so its geometric
generic fibre is integral.  Removing the exceptional divisor
does not alter that assertion.  For the local-system consequence,
apply the proved assertion to a connected component of its finite
frame cover, which is a torsor for its full monodromy image.
\end{proof}

\begin{corollary}\label{cor:restriction-independent-monodromy}
For $e\geq n+1$, fix a hyperplane $H\subset\mathbf P^n$.
The local system $V_{n,e}$ retains its full geometric monodromy
when the restriction of the equation to $H$ is prescribed to be
geometric generic.  On the parameter space where both the
equation and its restriction have invertible Hasse--Witt maps,
the joint monodromy image on $V_{n,e}$ and $V_{n-1,e}$ is the
product of their two full images.
\end{corollary}

\begin{proof}
Restriction of equations is a linear projection.  Its center
consists of equations divisible by the fixed equation of $H$.
It meets $W_{n,e}$: take $H$ and $e-1$ other general-position
hyperplanes.  Lemma~\ref{lem:linear-projection-monodromy} shows
that the joint image contains the first full group times $1$.
Its projection to the second group is surjective, since the
restriction map has geometrically integral generic fibre, an
open of an affine space.  The joint image is therefore the
product.  Removing further nonempty opens does not change
these finite images.
\end{proof}

\subsubsection{The arrangement module and independent quadratic reflections}

Take a general-position union of $d$ hyperplanes as base point.
Put $P=\mathbf F_p^d$, $u=\sum_i e_i$, $W=P/\langle u\rangle$,
$U=\ker(\sigma:P\to\mathbf F_p)$ with $\sigma(x)=\sum_i x_i$,
and $D=\operatorname{im}(U\to W)$.  The Cech resolution of the
union, whose strata are projective spaces, identifies its top
coherent cohomology and its Lang fibre with
\begin{equation}\label{eq:pn-arrangement-module}
 V_{n,d}=\bigwedge^nP/(u\wedge\bigwedge^{n-1}P)
       =\bigwedge^nW.
\end{equation}
The connected ordered-hyperplane cover of the unordered
arrangement locus gives the full permutation action of $S_d$
on this space.  Put $A=\bigwedge^nD$.

\begin{lemma}\label{lem:odd-general-hook-module}
Assume $1\leq n\leq d-2$.  If $p\nmid d$, then
$\bigwedge^nW$ is simple as an $S_d$-module.  If $p\mid d$, then
$A=\bigwedge^nD$ is its unique proper nonzero submodule; both
$A$ and the quotient $\bigwedge^{n-1}D$ are simple.
\end{lemma}

\begin{proof}
First suppose $2n\leq d$.  For $n$ pairwise disjoint transpositions
$\tau_j=(a_jb_j)$, the group-algebra operator
$\prod_j(1-\tau_j)$ has rank one on $\bigwedge^nW$.  Its image
is the wedge of the $n$ differences $e_{a_j}-e_{b_j}$, and its
functional is the wedge of the corresponding coordinate
differences.  This follows by simultaneously decomposing $W$
into the $n$ negative eigenlines and their common positive
eigenspace; $2$ is invertible.

These matching wedges span $\bigwedge^nU$.  To check this without
a modular representation-theory citation, use the difference
basis $f_i=e_i-e_o$ and the identity, for four distinct indices,
\[
 2f_a\wedge f_b
 =(e_a-e_b)\wedge(e_c-e_o)
  +(e_a-e_c)\wedge(e_b-e_o)
  +(e_a-e_o)\wedge(e_b-e_c).
\]
Beginning with a wedge of $n$ differences sharing $o$, repeatedly
introduce one unused index $c$.  Each application leaves one
fewer difference sharing $o$ and produces one disjoint pair.
There are $n-1$ available new indices because $d\geq2n$.
The resulting wedges are perfect matchings on $2n$ indices.
The same proof in the dual permutation space shows that the
matching functionals span $\bigwedge^nW^*$.
Thus every nonzero submodule contains a matching image and hence
all of $A$.  This proves simplicity of $A$, and proves simplicity
of the whole space when $p\nmid d$.

When $p\mid d$, the dot product gives a nondegenerate invariant
form on $D=U/\langle u\rangle$.  The same matching argument for
$n-1$ shows that $\bigwedge^{n-1}D$ is simple: matching images
span it, and their functionals span its dual under this form.
Contraction with the sum functional gives an exact sequence
\[
 0\longrightarrow\bigwedge^nD\longrightarrow\bigwedge^nW
 \longrightarrow\bigwedge^{n-1}D\longrightarrow0.
 \tag{*}
\]
Every nonzero submodule contains its kernel, and a submodule
which also surjects to the simple quotient is the whole space.
This proves the unique-submodule assertion in the stated range.

For $2n>d$, put $r=d-1-n$, so $1\leq r<d/2$.
The natural exterior pairing gives
$(\bigwedge^nW)^*\simeq\bigwedge^rW\otimes\operatorname{sgn}$.
If $p\nmid d$, the already proved simplicity for $r$ gives the
claim.  If $p\mid d$, duality preserves the property of having
exactly one proper nonzero submodule.  The displayed subspace
$\bigwedge^nD$ is nonzero and proper, so it must be that unique
submodule.  Its simplicity and that of the quotient follow
from the corresponding two simple factors for $r$.
\end{proof}

\begin{lemma}\label{lem:quadric-reflections}
For $n\geq2$ and $d\geq n+2$, the geometric group on
$V_{n,d}$ contains commuting rank-one reflections
$R_I$, indexed by $(n-1)$-subsets
$I\subset\{3,\ldots,d\}$.  At the arrangement base point,
their negative eigenlines are
\[
 \left\langle (e_1-e_2)\wedge\bigwedge_{i\in I}e_i\right\rangle,
\]
and their product is the permutation transposition $(12)$.
Here a rank-one reflection has eigenvalue $-1$ once and
eigenvalue $1$ on a codimension-one subspace.
\end{lemma}

\begin{proof}
Fix the last $d-2$ hyperplanes and replace the first two by a
variable quadric $Q$.  Require every positive-dimensional
quadric stratum to be connected with no higher coherent
cohomology, and every intersection with a line
$L_I=\bigcap_{i\in I}H_i$ to consist of two distinct points.
This is a nonempty open containing a smooth general quadric
and the reducible quadric $H_1H_2$.  A positive-dimensional
quadric hypersurface has $H^0(\mathcal O)=k$ and
$H^{>0}(\mathcal O)=0$, so the Cech complex computes all the
cohomology from its $H^0$ terms.  Only its top term has
nonconstant point sets: the two points on each $L_I$.
Swapping one pair acts on cohomology as a rank-one reflection,
since every coboundary has equal values on that pair.

All these swaps occur independently.  The restriction from
quadratic forms to $L_I\simeq\mathbf P^1$ is surjective.  The
discriminant of a binary quadratic is an irreducible polynomial
in its three coefficients in odd characteristic.  Its pullback
to the space of all quadratic forms is still irreducible, and
the discriminants for distinct lines are nonassociate.  To see
the last assertion, choose a point of $L_I\setminus L_J$ and
two hyperplanes through that point which meet $L_J$ in distinct
points.  Their product is tangent to $L_I$ and transverse to
$L_J$.  Their classes in the rational function field
modulo squares are therefore independent, by valuations at
these distinct prime divisors.  The ordered-pair cover has
Galois group $(\mathbf Z/2)^{\binom{d-2}{n-1}}$.
Removing the indicated open conditions leaves it connected.

The pair differences at $Q=H_1H_2$ are precisely the displayed
cochains, using the refinement of the cover with component $Q$
by its two hyperplanes.  Interchanging these two hyperplanes
interchanges every pair and fixes every other Cech term.
The transposition $(12)$ is thus the product of all the $R_I$.
Every fibre here has bijective Frobenius on the full top
coherent cohomology, because all terms are finite sets with
ordinary constant Frobenius.  The entire comparison is inside
$W_{n,d}$.
\end{proof}

\begin{proposition}\label{prop:pn-odd-irreducible}
For $n\geq2$ and $d\geq n+2$, the geometric monodromy group
$G_{n,d}$ is irreducible.
\end{proposition}

\begin{proof}
Only $p\mid d$ needs consideration, by
Lemma~\ref{lem:odd-general-hook-module}.  In that case any
proper invariant subspace must be $A=\bigwedge^nD$.
The negative eigenline of every $R_I$ lies outside $A$:
contraction with $\sigma$ sends its displayed generator to
\[
 -(e_1-e_2)\wedge\iota_\sigma
                  \left(\bigwedge_{i\in I}e_i\right)\ne0.
 \tag{**}
\]
Indeed the involved $n$ vectors remain linearly independent
modulo $u$: at least one index is outside $\{1,2\}\cup I$.
If $A$ were invariant, each $R_I$ would consequently fix $A$
pointwise.  Their product $(12)$ would do so too.  But on $D$
this transposition has a one-dimensional negative eigenspace
and a positive eigenspace of dimension $d-3$; on
$\bigwedge^nD$ its negative eigenspace has dimension
$\binom{d-3}{n-1}>0$.  This contradicts pointwise fixation.
\end{proof}


\subsubsection{The degree $n+2$ seed}

We use the following elementary observation about the Cech
computations, also in the determinant calculation below.  A family
consisting of coordinate hyperplanes and one further hypersurface
has the usual exact alternating resolution by its intersections
provided the latter equation restricts to a nonzero equation on
every coordinate stratum.  Locally its vanishing equation and any
of the coordinate equations form a regular sequence.  The
alternating union complex for a regular sequence is exact, by
induction using $(f)\cap(g)=(fg)$ and its versions modulo the
other equations.  Thus smoothness of the further hypersurface is
unnecessary.  Higher coherent cohomology of any hypersurface
stratum is given by its usual projective-space exact sequence.

\begin{proposition}\label{prop:pn-odd-seed-sl}
For $n\geq3$ and $d=n+2$, the group $G_{n,d}$ contains
$\mathrm{SL}_{n+1}(\mathbf F_p)$.
\end{proposition}

\begin{proof}
Let $N$ be the normal subgroup of $G_{n,d}$ generated by all its
transvections.  Replace any four of the arrangement hyperplanes,
indexed by $I$, by a varying quartic $Q$, and fix the other
$n-2$ hyperplanes.  The Cech spectral sequence has precisely
the following positive-cohomology terms contributing to total
degree $n-1$:
the quartic plane curve obtained by intersecting all the fixed
hyperplanes, of rank three, and the $n-2$ quartic surface strata,
each of rank one.  Higher-dimensional quartic strata have no
positive coherent cohomology.  The $H^0$ row has no positive
cohomology because the nerve is a full simplex.  There are no
nonzero incoming or outgoing differentials for the displayed
terms.  We consequently obtain an exact sequence of Frobenius
modules, and of Lang local systems on the invertibility open,
\[
 0\longrightarrow Q_I\longrightarrow V_{n,n+2}
      \longrightarrow T_I\longrightarrow0,
 \qquad \dim Q_I=3,\quad\dim T_I=n-2.
 \tag{1}
\]
Here $T_I$ is the direct sum of the rank-one surface local systems.
The open where all these terms are ordinary contains the original
arrangement.  Its image on $Q_I$ is $\mathrm{GL}_3(\mathbf F_p)$:
restriction of quartics to the fixed plane is a linear projection
with affine generic fibre, and the ordinary plane-quartic group
is already proved.  The inverse image of its connected frame
cover remains integral; imposing the other nonempty open
conditions preserves connectedness.  The image on $T_I$ is
abelian, being a subgroup of $(\mathbf F_p^\times)^{n-2}$.

Apply the quartic-boundary argument of the preceding plane
section to the dual of (1).  Goursat's lemma gives a normal
subgroup $B_0$ fixing $T_I^*$ pointwise and mapping onto
$\mathrm{SL}(Q_I^*)$.  Its kernel is
$Q_I\otimes T'$ for a subspace $T'\subseteq T_I^*$.
If the kernel is nonzero, it contains rank-one transvections.
If it is zero, Lemma~\ref{lem:sl3-natural-h1-odd} splits the
extension and again gives transvections.  Thus $N$ is nontrivial.
By Proposition~\ref{prop:pn-odd-irreducible} and
Lemma~\ref{lem:transvection-clifford-blocks}, $V_{n,n+2}^*$ is a
direct sum of canonical simple $N$-blocks permuted transitively
by $G_{n,n+2}$.

In the arrangement fibre there is the $S_d$-equivariant
identification
\[
 V_{n,n+2}^* \simeq W\otimes\operatorname{sgn},
 \qquad d=n+2.
 \tag{2}
\]
For a quartet $I$, let $U_I$ denote its three-dimensional
sum-zero subspace in $P$.  Since $4$ is invertible, it embeds
in $W$.  The restriction of (2) to its $S_4$ is the direct sum
of the unique simple module
$S_I=U_I\otimes\operatorname{sgn}$, of dimension three, and
$d-4$ copies of the sign character.  This is an actual
semisimple decomposition in every odd characteristic.
The quotient $Q_I^*$ in the dual of (1) is this unique
three-dimensional constituent.

The two cases in the proof of
Proposition~\ref{prop:odd-transvection-normal-irreducible} now
apply verbatim, with this specified decomposition: if the kernel
is zero, the commutator space of the split $\mathrm{SL}_3$ is
$S_I$ and belongs to one $N$-block; if it is nonzero, the
rank-one maps with a common quotient functional put $T'$ in one
block, and varying the functional makes that block surject to
$Q_I^*$.  The block is $S_4$-stable because $T'$ is stable, so
it contains the unique constituent $S_I$.

Quartets with three common indices give $S_I$'s with nonzero
intersection.  Their connected overlap graph puts all these
subspaces in one block $L$.  They span
$D\otimes\operatorname{sgn}$.  Every translate of $L$ by $S_d$
contains this same nonzero subspace and must equal $L$.
Thus $S_d$ also preserves the sum of the remaining blocks.
If that sum were nonzero, the degree-one case of
Lemma~\ref{lem:odd-general-hook-module}, with the sign twist,
would force it to contain $D\otimes\operatorname{sgn}$ too.
This contradicts the direct sum.  Therefore $N$ is irreducible.

Kantor--Seress, Theorem~3.6 \cite{KantorSeress2001}, applies:
over the prime field, powers of each transvection supply its
full transvection group.  It leaves
$N=\mathrm{SL}_{n+1}(\mathbf F_p)$ or a symplectic group.
The latter is excluded by
Lemma~\ref{lem:pn-reflection-excludes-symplectic} below.
\end{proof}

\begin{lemma}\label{lem:pn-reflection-excludes-symplectic}
If an irreducible normal subgroup $N\triangleleft G$ is the
natural symplectic group in dimension $h\geq3$ over an odd prime
field, then $G$ cannot contain a rank-one reflection.
\end{lemma}

\begin{proof}
The space of $N$-invariant alternating forms is one dimensional,
so normality puts $G$ in the symplectic similitude group.
Let $R$ be such a reflection, with positive eigenspace of
dimension $h-1$ and negative eigenline.  If its multiplier
were different from one, its positive eigenspace would be
totally isotropic, contradicting $h-1>h/2$.
Its multiplier is therefore one.  The negative eigenline is
then orthogonal to the positive eigenspace, since $p\ne2$,
and to itself because the form is alternating.  It is in the
radical of the form, a contradiction.
\end{proof}

\subsubsection{Induction in the degree}

\begin{proposition}\label{prop:pn-odd-degree-induction}
Suppose the groups for degree $d-1$ in $\mathbf P^n$ and
$\mathbf P^{n-1}$ contain their full special linear groups,
where $n\geq3$ and $d\geq n+3$.  Then
$G_{n,d}$ contains $\mathrm{SL}_{h_{n,d}}(\mathbf F_p)$.
If in addition the first of these two groups has full
determinant, then so does $G_{n,d}$.
\end{proposition}

\begin{proof}
Use the boundary $Y=A\cup H$, with $A$ a degree-$(d-1)$
hypersurface and $H$ a fixed hyperplane.  The usual exact
sequence of the union gives, on the simultaneous ordinary open,
\begin{equation}\label{eq:pn-two-piece-boundary}
 0\longrightarrow S\longrightarrow V_{n,d}
       \longrightarrow Q\longrightarrow0,
 \qquad
 S=V_{n-1,d-1},\quad Q=V_{n,d-1}.
\end{equation}
There are no adjacent positive coherent groups, since a
hypersurface in projective space has no intermediate coherent
cohomology.  Put
\[
 r=\dim Q=\binom{d-2}{n},\qquad
 s=\dim S=\binom{d-2}{n-1}.
\]
Both are at least three in the present range.
Corollary~\ref{cor:restriction-independent-monodromy} shows
that the graded image is the product of the full groups.
Take the inverse image $B_0$ of
$L=\mathrm{SL}(S)\times\mathrm{SL}(Q)$.
Its kernel $K$ is an additive $L$-submodule of
$M=\operatorname{Hom}(Q,S)$.  This outer tensor product of
the two absolutely simple natural modules is absolutely simple,
so $K=0$ or $K=M$.

If $K=M$, its rank-one elements are transvections of $V$.
Let $N$ be their full normal closure in $G_{n,d}$ and use
the canonical $N$-blocks, justified by irreducibility of $G$.
For every $w\in S$ and $\phi\in Q^*$, the operator
$1+w\otimes\phi$ is a transvection when both are nonzero.
Fixing $\phi$ puts all of $S$ in one block; fixing a nonzero
$w$ and varying $\phi$ makes that block surject to $Q$.
It contains $S$ and surjects to $V/S$, so it is all of $V$.
Thus $N$ is irreducible.

If $K=0$, the complement is conjugate to the split block
diagonal group $L$.  Here is the needed cohomology computation,
which requires no vanishing theorem for either natural module:
\[
 H^1(\mathrm{SL}(S)\times\mathrm{SL}(Q),S\otimes Q^*)=0.
\]
The first factor has no invariants in $S$.  In the low-degree
Hochschild--Serre sequence the remaining term is
$H^1(\mathrm{SL}(S),S)\otimes Q^*$, whose invariants under the
second factor are zero.  This proves the displayed vanishing
and hence the conjugacy of complements.
The split $L$ is generated by rank-one transvections on $V$,
so is contained in $N$.  Every canonical $N$-block is therefore
an $L$-submodule of the semisimple, multiplicity-free module
$S\oplus Q$.  If there is more than one block, they are exactly
these two summands.  Transitivity of $G$ on the blocks forces
$r=s$, equivalently $d=2n+1$.

Even this last possibility is impossible.  A rank-one reflection
cannot interchange two subspaces of common dimension $r>1$:
its difference from the identity would have rank at least $r$
on either one.  Thus every reflection of
Lemma~\ref{lem:quadric-reflections} preserves both blocks.
Their product $(12)$ does so too.  Conjugating within the
arrangement $S_d$ shows that every transposition, hence all
of $S_d$, preserves both blocks.  This contradicts
Lemma~\ref{lem:odd-general-hook-module}: a simple module, or a
module with a unique proper nonzero submodule, has no such
nontrivial direct sum decomposition.  We have proved that
$N$ is irreducible also when $K=0$.

The transvection classification again gives $N=\mathrm{SL}(V)$
or a symplectic group.  The latter is excluded by
Lemma~\ref{lem:pn-reflection-excludes-symplectic}, using the
quadric reflection already present in $G$.
Finally the full product assertion for the graded image
contains the old group on $Q$ times $1$ on $S$.  Its lifts
have the same determinant as their action on $Q$, because
the extension kernel is unipotent.  Full determinant for
that old group therefore implies full determinant for $G$.
\end{proof}

\subsubsection{Determinants in the two initial degrees}

\begin{lemma}\label{lem:elliptic-fixed-three-points}
There is a family of plane cubics over a strictly henselian
smooth trait over $k$ whose special fibre is smooth supersingular,
whose Hasse invariant has valuation one, and whose restriction
to a fixed line is the constant squarefree binary cubic
$g=x_0x_1(x_0-x_1)$.  There are also smooth ordinary plane
cubics with this same restriction.
\end{lemma}

\begin{proof}
Take the Legendre family at a simple supersingular zero of its
Hasse polynomial, as proved in the earlier elliptic calculation.
Choose a line transverse to the special fibre and identify it
with $\mathbf P^1$.  Transversality persists on the trait, and
the three intersection points form a finite etale scheme.
They split after strict henselization.  The unique projective
linear transformation on the line taking their chosen order
to $0,1,\infty$ varies regularly with the trait parameter.
Lift this transformation to the ambient projective plane and
rescale the cubic equation by a unit.  Its restriction to the
line is then exactly $g$.  Isomorphisms and unit changes of
cohomology trivialization multiply the Hasse invariant by a
unit, so its valuation is still one.  The same construction
with any ordinary elliptic plane cubic gives the last assertion.
\end{proof}

\begin{proposition}\label{prop:pn-seed-full-determinant}
For every $n\geq2$, the geometric determinant images in degrees
$n+1$ and $n+2$ are the full group $\mathbf F_p^\times$.
\end{proposition}

\begin{proof}
For degree $n+1$, fix $n-2$ coordinate hyperplanes and let their
intersection be a plane.  Lift the cubic family of
Lemma~\ref{lem:elliptic-fixed-three-points} from this plane to
$\mathbf P^n$, for instance by its same equation independent
of the other coordinates, and take its union with the fixed
hyperplanes.  Their intersections are proper.  All
positive-dimensional cubic strata of ambient dimension at
least three have no positive coherent cohomology; the only
nonzero term in total degree $n-1$ is the elliptic curve term.
The $H^0$ nerve is a full simplex and has no higher cohomology.
Thus the rank-one total Hasse invariant is a unit times the
elliptic one and has valuation exactly one.
For $n=2$ this is simply the elliptic family itself.

For degree $n+2$, use coordinates
$x_0,x_1,y_1,\ldots,y_{n-1}$ and hyperplanes $H_i=(y_i=0)$.
Their common line is $\ell=\mathbf P(x_0,x_1)$.
Let $f_1(x_0,x_1,y_1)$ be the family in the lemma, and for
$i>1$ let $f_i(x_0,x_1,y_i)$ be a fixed ordinary plane cubic,
all restricting to $g$ on $\ell$.  Define explicitly
\[
 q=g+\sum_{i=1}^{n-1}(f_i-g),\qquad
 Y=(q\,y_1\cdots y_{n-1}=0)\subset\mathbf P^n.
 \tag{3}
\]
Restriction to the plane where all $y_j$, $j\ne i$, vanish
is exactly $f_i$.  Every coordinate stratum has a nonzero
restricted cubic equation, because its restriction to $\ell$
is $g$.  The alternating union resolution described above
is therefore exact, even if the higher-dimensional cubic
strata happen to be singular.

Those higher-dimensional cubic strata have no positive
coherent cohomology, and every positive-dimensional stratum
has $H^0(\mathcal O)=k$.  The common zero-dimensional stratum
consists of the three constant, distinct roots of $g$.
The $H^0$ row is the full-simplex complex except that its last
copy of $k$ is replaced by $k^3$; it contributes an ordinary
rank-two piece in degree $n-1$.  In the positive-cohomology
rows the $n-1$ elliptic curves contribute their rank-one
cohomology, again in total degree $n-1$.  There are no
possible nonzero differentials.  This gives a Frobenius-stable
filtration on the total cohomology with ordinary graph piece
of rank two and the $n-1$ elliptic pieces.  Its total rank is
$2+(n-1)=n+1=\binom{n+1}{n}$.
The Hasse determinant is a unit times the product of their
Hasse invariants.  Exactly $f_1$ has valuation one and all the
other factors are units, so the total determinant has
valuation one.

In both constructions the punctured trait lies in $W_{n,d}$.
The determinant Lang torsor is
$z^{p-1}=h^{-1}$, as in \eqref{eq:determinant-kummer}.
A valuation-one $h$ gives tame geometric inertia of order
$p-1$.  This proves the assertion.  The special fibre is
allowed to have smaller Hasse--Witt rank: only its valuation
is used, and no stable local system is extended across it.
\end{proof}

\begin{theorem}\label{thm:pn-odd-full-groups}
Let $p$ be odd, let $k$ be any perfect field of characteristic
$p$, and let $n\geq2$.  On the ordinary smooth degree-$d$
hypersurface locus in $\mathbf P^n_k$, the geometric and
arithmetic monodromy images on
$H^{n-1}_{\mathrm{et}}(Y_{\bar\eta},\mathbf F_p)$ are both
\[
 \mathrm{GL}_{\binom{d-1}{n}}(\mathbf F_p).
\]
For $d\leq n$ the vector space is zero and the image is trivial.
For $d=n+1$ it has rank one and full determinant.
For every $d\geq n+2$ it is irreducible with full determinant.
\end{theorem}

\begin{proof}
Dimension two is the already proved plane theorem.
Induct on $n\geq3$.  The zero-rank degrees need no argument,
and Proposition~\ref{prop:pn-seed-full-determinant} proves the
rank-one degree $n+1$.  In degree $n+2$,
Proposition~\ref{prop:pn-odd-seed-sl} and the determinant
proposition give the full general linear group.
Now induct on $d\geq n+3$ using
Proposition~\ref{prop:pn-odd-degree-induction}.
The two smaller groups are known by the dimension and degree
inductions, and the old determinant propagates as proved there.
This proves the geometric image over $\overline{\mathbf F}_p$.
The connected finite frame cover remains geometrically
connected under algebraically closed extension; since the
family is defined over $\mathbf F_p$, the result holds over
the algebraic closure of any perfect $k$.  The arithmetic
image contains the full geometric general linear group, so
is equal to it.
\end{proof}

\begin{corollary}\label{cor:pn-odd-all-residues-sharp}
Let $k$ be finite of odd characteristic $p$ and $n\geq2$.
For every $d\geq n+2$, every nonempty open in the smooth
degree-$d$ hypersurface locus of $\mathbf P^n_k$ contains,
over suitable finite extensions of $k$, hypersurfaces with
each of the $p$ point-count residues.
For $d\leq n$ every smooth hypersurface has residue $1$.
For $d=n+1$ the ordinary open misses residue $1$.
Thus $d=n+2$ is the exact threshold for the assertion involving
every nonempty open.
\end{corollary}

\begin{proof}
The finite-field coherent trace formula gives
\[
 \#Y(k')\equiv 1+(-1)^{n-1}
       \operatorname{Tr}(F_{k'}\mid H^{n-1}(Y,\mathcal O_Y)_{\rm ss})
       \pmod p,
\]
because a projective hypersurface has no other positive
coherent cohomology.  In the zero-rank range the trace is zero.
In rank one on the ordinary open it is nonzero.
In rank $h\geq2$, traces of invertible matrices take every
value in $\mathbf F_p$: the block
$\begin{pmatrix}0&-1\\1&a\end{pmatrix}$ has determinant one
and trace $a$, and adjoining an identity block allows any
specified total trace.  By the theorem the finite geometric
and arithmetic image is the full group.  The finite
Chebotarev theorem already cited in the plane calculation
therefore realizes each such conjugacy class in every
nonempty open after suitable finite extension of $k$.
The intersection of that open with the ordinary locus is
nonempty because the latter is dense.
\end{proof}
